% SPDX-License-Identifier: CC-BY-4.0
\section{\sediment{}}
\label{sec:sediment}

\subsection{Algorithm}

\sediment{} models each feature as an independent Gaussian under the
clean null. The score on a new observation $x$ is the sum of squared
standardized residuals,
\begin{equation}
  \score(x) = \sum_{j=1}^{d} \frac{(x_j - \mu_j)^2}{(\sigma_j + \epsilon)^2},
  \label{eq:score}
\end{equation}
which is the Mahalanobis distance under a diagonal-covariance
assumption. The estimator has two parameters: a trim fraction
$\trim \in [0, 1)$ and a small $\epsilon > 0$ that floors the
variance.

For each feature $j$, given a learning window
$\{x_1^{(j)}, \dots, x_N^{(j)}\}$ sorted ascending, the per-tail
trim count is $k = \lfloor \trim N / 2 + 1/2 \rfloor$. The trimmed
mean is the mean of the surviving $N - 2k$ samples,
\begin{equation}
  \mu_j = \frac{1}{N - 2k} \sum_{i=k+1}^{N-k} x_{(i)}^{(j)}.
\end{equation}
The trimmed standard deviation is the unbiased sample standard
deviation of the surviving samples, multiplied by a closed-form
Gaussian-truncation correction $c(\trim)$ derived below,
\begin{equation}
  \sigma_j = c(\trim) \cdot \sqrt{\frac{1}{N - 2k - 1}
                       \sum_{i=k+1}^{N-k} (x_{(i)}^{(j)} - \mu_j)^2}.
\end{equation}

\subsection{Gaussian-truncation correction}

For an $\mathcal{N}(0, s^2)$ sample symmetrically trimmed at
$\trim/2$ per tail, the trimmed-sample variance underestimates $s^2$
because the trimmed tails carried disproportionate variance. Let
$\Phi^{-1}$ and $\phi$ denote the standard normal quantile and
density functions. The trim cutoff is
\begin{equation}
  z = \Phi^{-1}\!\left(1 - \frac{\trim}{2}\right).
\end{equation}
The variance of an $\mathcal{N}(0, 1)$ variable conditioned on
$|Z| \leq z$ is
\begin{equation}
  \mathrm{Var}(Z \mid |Z| \leq z) = 1 - \frac{2 z \phi(z)}{1 - \trim},
  \label{eq:truncated-variance}
\end{equation}
so the multiplicative correction that recovers an unbiased estimate
of $s$ from the trimmed sample standard deviation is
\begin{equation}
  c(\trim) = \frac{1}{\sqrt{1 - 2 z \phi(z) / (1 - \trim)}}.
\end{equation}

\subsection{Robustness guarantee}

The classical breakdown point of the symmetric $\trim$-trimmed mean
is $\trim/2$~\cite{huber1964robust}: an adversary controlling fewer
than $\trim N / 2$ samples per feature cannot move the trimmed mean
an arbitrary amount. We adopt this as our defender's contract.

\begin{proposition}[\sediment{} robustness, informal]\label{prop:robust}
Fix a per-feature trim fraction \trim{} and an adversary budget
$\budget \leq \trim/2$. Let $\mu^{\circ}_j, \sigma^{\circ}_j$ denote
the population mean and standard deviation of feature $j$ under the
clean data-generating distribution, and let
$\hat{\mu}_j, \hat{\sigma}_j$ denote the moments fitted by
\sediment{} on a poisoned learning window of size $N$. Then for each
feature $j$, $\hat{\mu}_j \to \mu^{\circ}_j$ and
$\hat{\sigma}_j \to \sigma^{\circ}_j$ in probability as
$N \to \infty$, regardless of the adversary's choice of poisoning
placement.
\end{proposition}

The argument has two ingredients. First, the trimmed mean's
breakdown point of $\trim/2$~\cite{huber1964robust} bounds the
influence of any $\budget N$ adversarial samples on $\hat{\mu}_j$
when $\budget \leq \trim/2$, because the adversarial samples cannot
all sit in the surviving region simultaneously. Second, the
closed-form correction $c(\trim)$ restores the trimmed standard
deviation to an unbiased estimate of $\sigma^{\circ}_j$ on the
surviving $(1 - \trim) N$ benign samples.

The bound is tight in the worst case. At $\budget = \trim/2$, an
adversary can in principle place every adversarial sample at the
defender's trim boundary, surviving the trim and dragging
$\hat{\mu}_j$. \cref{sec:eval-white-box} reports a grid-search
white-box adversary that attempts exactly this. Empirically the
optimum reduces to the at-target injection, which gives evidence
that the worst case is hard to construct in practice. We return to
residual risks in \cref{sec:discussion}.

\subsection{Complexity}

Per feature, fitting requires sorting the $N$ samples
($\mathcal{O}(N \log N)$) and computing the trimmed moments
($\mathcal{O}(N)$). Memory is $\mathcal{O}(d)$ for the fitted
moments per process. Scoring is $\mathcal{O}(d)$ per observation.
